\documentclass[a4paper,12pt]{article}
\usepackage[left=2cm,right=2cm,top=2cm,bottom=2.5cm]{geometry}
\usepackage[utf8]{inputenc}
\usepackage[T1]{fontenc}
\usepackage[french]{babel}

\usepackage{iftex}
\ifPDFTeX
  \usepackage{fourier}       % pdfLaTeX (polices Type1)
\else
  \usepackage{fourier-otf}   % LuaLaTeX ou XeLaTeX (polices OpenType)
\fi

\usepackage{tocloft}
\renewcommand{\cftsecleader}{\cftdotfill{\cftdotsep}}  % points de conduite dans la table des matières

\usepackage{hyperref}
\hypersetup{colorlinks=true, linkcolor=blue, urlcolor=blue}
\usepackage{xcolor}
\usepackage{tabularx}
\usepackage{array}
\usepackage{booktabs}
\usepackage{pgfplots}
\pgfplotsset{compat=1.18}
\usetikzlibrary{babel,patterns}
\pgfplotsset{
  /pgf/number format/use comma,
  /pgf/number format/1000 sep={\,},
  planche/.style={
    width=14.5cm, height=10cm,
    grid=both,
    grid style={gray!30},
    minor grid style={gray!15},
    tick align=outside,
    every axis/.append style={font=\small},
    label style={font=\small},
  },
}

\newcommand{\degmin}[2]{$#1^{\circ}\,#2'$}
\renewcommand{\deg}[1]{$#1^{\circ}$}  % remplace l'opérateur mathématique deg, inutilisé ici

\title{Planches du manuel de vol\\
\Large Alouette II, turbine Artouste II C}
\author{}
\date{Reconstitution du 11 août 2026}

\begin{document}
\maketitle

\noindent\fbox{\parbox{\dimexpr\linewidth-2\fboxsep-2\fboxrule}{%
\textbf{Avertissement.} Ce document est une reconstitution typographique et
graphique des planches du manuel de vol allemand de l'Alouette II (turbine
Artouste II C), d'après le scan diffusé par FSHeli.ch ("for Flight Simulator
use only"), archivé sous
\texttt{artouste-sources-tierces/docs-se3130/fsheli-alouette2-reference.pdf}.
Les courbes ont été numérisées sur le scan le 11/08/2026 ; l'incertitude de
lecture est indiquée sous chaque planche. Document destiné à la simulation
uniquement, sans aucune valeur opérationnelle.}}

\tableofcontents

\section{Vitesse de montée à 90 km/h (Abb. 3-6)}

Conditions de la planche : toutes configurations (patins ou flotteurs), pas
collectif \deg{14}, régime 34\,000 tr/min, vitesse indiquée 90 km/h (50 kt).
La planche d'origine se lit en deux temps : conversion de l'altitude-pression
en altitude-densité selon la température, puis lecture de la vitesse de
montée à l'altitude-densité obtenue.

\begin{center}
\begin{tikzpicture}
\begin{axis}[planche, height=8cm,
  xlabel={Altitude-densité (m)}, ylabel={Altitude-pression (m)},
  xmin=0, xmax=4600, ymin=0, ymax=4500,
  minor x tick num=1, minor y tick num=1]
\addplot[gray, thin] coordinates {(-1160,0) (3331,4500)}
  node[pos=0.93, above left, black, font=\scriptsize, rotate=42]{$-30$};
\addplot[gray, thin] coordinates {(-756,0) (3740,4500)}
  node[pos=0.97, above left, black, font=\scriptsize, rotate=42]{$-20$};
\addplot[gray, thin] coordinates {(-369,0) (4130,4500)}
  node[pos=0.99, above left, black, font=\scriptsize, rotate=42]{$-10$};
\addplot[black, thick] coordinates {(0,0) (4500,4500)}
  node[pos=0.93, above left, font=\scriptsize, rotate=42]{ISA};
\addplot[gray, thin] coordinates {(354,0) (4853,4500)}
  node[pos=0.85, below right, black, font=\scriptsize, rotate=42]{$+10$};
\addplot[gray, thin] coordinates {(693,0) (5190,4500)}
  node[pos=0.78, below right, black, font=\scriptsize, rotate=42]{$+20$};
\addplot[gray, thin] coordinates {(1020,0) (5513,4500)}
  node[pos=0.71, below right, black, font=\scriptsize, rotate=42]{$+30$};
\end{axis}
\end{tikzpicture}

\begin{tikzpicture}
\begin{axis}[planche,
  xlabel={Altitude-densité (m)}, ylabel={Vitesse de montée (m/s)},
  xmin=0, xmax=4600, ymin=0, ymax=9,
  minor x tick num=1, minor y tick num=1,
  legend style={font=\scriptsize, at={(0.98,0.98)}, anchor=north east},
  legend cell align=left]
\addplot[black, thick, smooth] coordinates
  {(0,8.10) (1000,7.00) (2000,5.95) (2500,5.55) (3000,5.10) (3500,4.30) (4000,3.85) (4500,3.25)};
\addlegendentry{1100 kg}
\addplot[black, smooth] coordinates
  {(0,7.20) (1000,6.15) (2000,5.15) (2500,4.64) (3000,4.00) (3500,3.45) (4000,2.85) (4500,2.10)};
\addlegendentry{1200 kg}
\addplot[black, smooth] coordinates
  {(0,6.40) (1000,5.35) (2000,4.35) (2500,3.67) (3000,3.10) (3500,2.32) (4000,1.95) (4450,0.00)};
\addlegendentry{1300 kg}
\addplot[black, smooth] coordinates
  {(0,5.55) (1000,4.62) (2000,3.55) (2500,2.76) (3000,2.08) (3500,1.17) (3950,0.00)};
\addlegendentry{1400 kg}
\addplot[black, smooth] coordinates
  {(0,4.89) (1000,3.87) (2000,2.78) (2500,1.70) (3000,0.55) (3080,0.00)};
\addlegendentry{1500 kg}
\addplot[black, dashed, smooth] coordinates
  {(0,4.15) (1000,3.10) (2000,1.95) (2300,0.00)};
\addlegendentry{1600 kg}
\end{axis}
\end{tikzpicture}
\end{center}

Lecture numérisée à $\pm 0{,}15$ m/s. La pente commune est d'environ
$-1{,}05$ m/s par 1\,000 m d'altitude-densité, jusqu'à un effondrement final
environ 2 m/s au-dessus du plafond de montée. Plafonds lus : environ
2\,300 m à 1600 kg, 3\,050 m à 1500 kg, 3\,950 m à 1400 kg, 4\,450 m à
1300 kg ; hors planche au-delà. L'axe des montées de la planche d'origine
est orienté vers le bas ; il est redressé ici.

\section{Carte moteur (Abb. 0-8)}

Puissance en fonction de la température de tuyère et de la consommation
horaire, en grandeurs corrigées (réduites aux conditions standard). Réseau
partiel : trois iso-températures et trois iso-consommations sur les réseaux
complets de la planche (420 à 510 \deg{}C par pas de 10, 140 à 190 kg/h par
pas de 5).

\begin{center}
\begin{tikzpicture}
\begin{axis}[planche, height=11cm,
  xlabel={$\dfrac{n}{1000}\sqrt{\dfrac{288}{t_0}}$},
  ylabel={$P\,\dfrac{760}{P_a}\sqrt{\dfrac{288}{t_0}}$\quad (PS)},
  xmin=31, xmax=35.4, ymin=270, ymax=500,
  minor x tick num=1, minor y tick num=1]
\addplot[black, thick, smooth] coordinates
  {(31.5,412) (32,431) (32.5,448) (33,461) (33.5,470) (34,476) (34.4,478) (34.7,476)}
  node[pos=0.55, above left, font=\scriptsize]{510};
\addplot[black, thick, smooth] coordinates
  {(31.7,375) (32,388) (32.5,404) (33,415) (33.5,424) (34,428) (34.3,429) (34.7,419)}
  node[pos=0.55, above left, font=\scriptsize]{480};
\addplot[black, thick, smooth] coordinates
  {(31.9,341) (32.5,356) (33,364) (33.5,371) (34,375) (34.1,375) (34.5,370) (34.8,345)}
  node[pos=0.5, above left, font=\scriptsize]{450};
\addplot[black, dashed, smooth] coordinates
  {(33.2,480) (33.5,478) (34,474) (34.5,467) (34.8,459)}
  node[pos=0.75, above right, font=\scriptsize]{190};
\addplot[black, dashed, smooth] coordinates
  {(31.7,410) (32,413) (32.5,415) (33,412) (33.5,407) (34,400) (34.5,385) (34.8,373)}
  node[pos=0.85, above right, font=\scriptsize]{170};
\addplot[black, dashed, smooth] coordinates
  {(31.6,326) (32,333) (32.5,336) (33,331) (33.5,322) (34,315) (34.5,298)}
  node[pos=0.85, above right, font=\scriptsize]{150};
\end{axis}
\end{tikzpicture}
\end{center}

\noindent Grandeurs : $P$ puissance (PS), $n$ régime du moteur (tr/min),
$P_a$ pression atmosphérique (mm Hg), $t_0$ température extérieure absolue
(K), $C_h$ consommation (kg/h), $t_4$ température de tuyère (\deg{}C).
Traits pleins : $t_4$ corrigée $(t_4+273)\,\frac{288}{t_0}-273$. Tirets :
consommation corrigée $C_h\,\frac{760}{P_a}\sqrt{\frac{288}{t_0}}$.
Lecture à $\pm 4$ PS.

\section{Puissance maximale de décollage (Abb. 0-9 et 0-10)}

\subsection{Au niveau de la mer, selon la température (Abb. 0-9)}

\begin{center}
\begin{tikzpicture}
\begin{axis}[planche,
  xlabel={Température extérieure (\deg{}C)},
  ylabel={Puissance réelle (PS)},
  xmin=-35, xmax=55, ymin=295, ymax=465,
  minor x tick num=1, minor y tick num=1]
\addplot[black, dashed, thick] coordinates {(18,455) (50,372)}
  node[pos=0.35, right=2pt, font=\scriptsize]{$t_4$ max, 34\,000};
\addplot[black, dashed] coordinates {(20,444) (50,364)}
  node[pos=0.72, below left, font=\scriptsize]{$t_4$ max, 33\,000};
\addplot[black, densely dashed, thick] coordinates {(-30,391) (0,404) (10,406) (34,406)}
  node[pos=0.45, above, font=\scriptsize]{débit carburant, 33\,000};
\addplot[black, densely dashed] coordinates {(-20,381) (0,398) (10,401) (39,402)}
  node[pos=0.62, below, font=\scriptsize]{débit carburant, 34\,000};
\addplot[black, thick] coordinates {(-30,390) (50,313)}
  node[pos=0.62, above right, font=\scriptsize]{butée de pas, 34\,000};
\addplot[black] coordinates {(-30,377) (50,316)}
  node[pos=0.35, below left, font=\scriptsize]{butée de pas, 33\,000};
\end{axis}
\end{tikzpicture}
\end{center}

La butée de pas est \degmin{14}{30} à $V_{\text{anz}} = 90$ km/h (50 kt) ou
\deg{14} à $V_{\text{anz}} = 0$. En dessous d'environ $+35$ \deg{}C, la
puissance utilisable au décollage est bornée par la butée de pas, très en
dessous des limites moteur. Lecture à $\pm 5$ PS.

\subsection{En atmosphère normale, selon l'altitude (Abb. 0-10)}

Régime 34\,000 tr/min.

\begin{center}
\begin{tikzpicture}
\begin{axis}[planche, height=9cm,
  xlabel={Puissance (PS)}, ylabel={Altitude-pression (m)},
  xmin=220, xmax=490, ymin=0, ymax=3300,
  minor x tick num=1, minor y tick num=1]
\addplot[black, thick] coordinates {(357,0) (228,3200)}
  node[pos=0.45, left=2pt, font=\scriptsize, align=center]{butée\\de pas};
\addplot[black, densely dashed] coordinates {(400,0) (443,2400)}
  node[pos=0.5, right=2pt, font=\scriptsize, align=center]{débit carburant\\170 kg/h};
\addplot[black, dashed, thick] coordinates {(475,0) (398,3200)}
  node[pos=0.75, right=2pt, font=\scriptsize]{$t_4$ max};
\end{axis}
\end{tikzpicture}
\end{center}

La limite moteur passe du débit carburant à la température $t_4$ vers
1\,800 m ; la butée de pas (\degmin{14}{30} à 90 km/h ou \deg{14} à
0 km/h) reste la borne réelle sur toute la plage. Lecture à $\pm 5$ PS.

\section{Butées du pas collectif en utilisation normale (1.6a)}

Ces limites respectées, la température $t_4$ maximale n'est jamais
dépassée.

\begin{center}
\small
\begin{tabular}{@{}llc@{}}
\toprule
Phase de vol & Situation & Pas maximal \\
\midrule
Stationnaire & & \deg{14} \\
\addlinespace
Décollage normal (terrain dégagé) & en stationnaire & \deg{14} \\
 & passage en translation & \degmin{14}{30} \\
\addlinespace
Décollage à forte pente & en stationnaire & \degmin{13}{30} \\
 & survol des obstacles & \degmin{14}{30} \\
\addlinespace
Montée ($V_i = 90$ km/h, 33\,000 tr/min) & & \deg{14} \\
\addlinespace
Croisière & 33\,000 tr/min & \degmin{14}{30} \\
 & 34\,000 tr/min & \deg{14} \\
\bottomrule
\end{tabular}
\end{center}

\paragraph{Huile froide.} En dessous de $-10$ \deg{}C avec l'huile O-138,
faire un point fixe au sol et attendre que l'aiguille de température
d'huile monte. Passer en translation sans augmenter le pas, ou rester en
stationnaire, jusqu'à 20 \deg{}C d'huile avec l'O-138, 0 \deg{}C avec
l'O-135.

\paragraph{Avertissements.} Ne jamais augmenter le pas brutalement. En
utilisation normale, ne jamais franchir la butée élastique. En cas
d'urgence, la butée élastique peut être dépassée jusqu'à \deg{15} sans
risque de pompage.

\section{Vitesses à ne pas dépasser (1.7.1)}

Vitesses limites $V_{\text{anz}}$ pour tous les états de vol, patins ou
flotteurs. La vitesse indiquée est égale à la vitesse corrigée. Valeurs en
km/h, nœuds entre parenthèses.

\begin{center}
\small
\begin{tabular}{@{}rrrrrrrr@{}}
\toprule
\multicolumn{2}{c}{Masse au décollage} & \multicolumn{6}{c}{Altitude (m)} \\
\cmidrule(lr){1-2}\cmidrule(l){3-8}
kg & lb & 0 & 1000 & 2000 & 3000 & 4000 & 4500 \\
\midrule
1600 & 3500 & 185 (100) & 165 (90) & & & & \\
1500 & 3300 & 195 (105) & 175 (95) & 155 (85) & & & \\
1400 & 3100 & 195 (105) & 185 (100) & 165 (90) & 145 (80) & & \\
1300 & 2900 & 195 (105) & 195 (105) & 175 (95) & 155 (85) & 135 (75) & \\
1200 & 2700 & 195 (105) & 195 (105) & 185 (100) & 165 (90) & 145 (80) & 135 (75) \\
1100 & 2400 & 195 (105) & 195 (105) & 195 (105) & 175 (95) & 155 (85) & 145 (80) \\
1000 & 2200 & 195 (105) & 195 (105) & 195 (105) & 185 (100) & 165 (90) & 155 (85) \\
\bottomrule
\end{tabular}
\end{center}

Approximations données par le manuel : retrancher 20 km/h (10 kt) par
1\,000 m de hauteur ; ajouter 10 km/h (5 kt) par 100 kg de masse en moins,
sans jamais dépasser 195 km/h (105 kt).

\section{Domaine de sécurité pour atterrissage en autorotation (1.8)}

Valable à 1600 kg, en vol face au vent. À masse plus faible, la zone à
éviter est plus petite : la courbe publiée sert de borne prudente.

\begin{center}
\begin{tikzpicture}
\begin{axis}[planche, height=11cm,
  xlabel={Vitesse indiquée (km/h)}, ylabel={Hauteur (m)},
  xmin=0, xmax=190, ymin=0, ymax=160,
  minor x tick num=1, minor y tick num=1]
\addplot[black, fill=gray!35] coordinates
  {(0,150) (30,136) (40,127) (50,117) (60,100) (65,90) (70,76) (72,63)
   (70,48) (65,37) (60,30) (50,21) (40,13) (30,8) (20,5) (10,3) (0,2)} -- cycle;
\addplot[black, fill=gray!35] coordinates
  {(52,0) (60,5) (70,9) (80,13) (90,17) (100,20) (120,22) (140,24) (185,25) (185,0)} -- cycle;
\node[font=\small\itshape, align=center] at (axis cs:33,70) {zone à\\éviter};
\node[font=\small\itshape, align=center] at (axis cs:130,80) {domaine de sécurité};
\node[font=\small\itshape, align=center] at (axis cs:125,10) {zone à éviter};
\end{axis}
\end{tikzpicture}
\end{center}

Numérisation à $\pm 3$ km/h et $\pm 3$ m. Le bec de la zone basse vitesse
est à 72 km/h et 63 m ; le survol au ras du sol (sous 2 m en
stationnaire) et le vol au-delà du bec sont sûrs. La zone grande vitesse
impose de tenir environ 25 m de hauteur au-delà de 140 km/h.

\section{Contrôle de la puissance (Abb. 2-23)}

Température de tuyère $t_4$ en fonction du pas collectif, à
34\,000 tr/min, pour le contrôle périodique de la puissance. La planche de
l'exemplaire scanné porte une correction individuelle de turbine de
$+15$ \deg{}C.

\begin{center}
\begin{tikzpicture}
\begin{axis}[planche,
  xlabel={Pas collectif (\deg{})}, ylabel={$t_4$ (\deg{}C)},
  xmin=12, xmax=14.35, ymin=340, ymax=525,
  minor x tick num=3, minor y tick num=1]
\addplot[black] coordinates {(12,455) (14,512)} node[right, font=\scriptsize]{$+45$};
\addplot[black] coordinates {(12,439) (14,495)} node[right, font=\scriptsize]{$+35$};
\addplot[black] coordinates {(12,422) (14,478)} node[right, font=\scriptsize]{$+25$};
\addplot[black] coordinates {(12,407) (14,461)} node[right, font=\scriptsize]{$+15$};
\addplot[black] coordinates {(12,391) (14,443)} node[right, font=\scriptsize]{$+5$};
\addplot[black] coordinates {(12,375) (14,426)} node[right, font=\scriptsize]{$-5$};
\addplot[black] coordinates {(12,358) (14,410)} node[right, font=\scriptsize]{$-15$};
\node[font=\scriptsize, rotate=90] at (axis cs:14.28,470) {température extérieure (\deg{}C)};
\end{axis}
\end{tikzpicture}
\end{center}

Pente lue : 26 à 28 \deg{}C de $t_4$ par degré de pas, environ
1,6 \deg{}C par degré de température extérieure. Lecture à $\pm 3$ \deg{}C.

\section{Contrôle du pas en stationnaire (Abb. 2-32)}

Pas collectif nécessaire au stationnaire à 1,5 m du sol, masse
$1500 \pm 20$ kg, 34\,000 tr/min, vent faible.

\begin{center}
\begin{tikzpicture}
\begin{axis}[planche, height=10.5cm,
  xlabel={Pas collectif (\deg{})}, ylabel={Altitude-pression (m)},
  xmin=12, xmax=14, ymin=0, ymax=2100,
  minor x tick num=3, minor y tick num=1,
  clip=true]
% iso-températures tous les 5 degrés C, de -20 à +40
\addplot[gray]  coordinates {(12.03,0) (13.43,2000)} node[pos=0.4, left, black, font=\scriptsize]{$-20$};
\addplot[gray!60] coordinates {(12.14,0) (13.54,2000)};
\addplot[gray]  coordinates {(12.25,0) (13.65,2000)} node[pos=0.5, left, black, font=\scriptsize]{$-10$};
\addplot[gray!60] coordinates {(12.36,0) (13.76,2000)};
\addplot[gray]  coordinates {(12.47,0) (13.87,2000)} node[pos=0.6, left, black, font=\scriptsize]{$0$};
\addplot[gray!60] coordinates {(12.58,0) (13.98,2000)};
\addplot[gray]  coordinates {(12.69,0) (14.09,2000)} node[pos=0.68, left, black, font=\scriptsize]{$+10$};
\addplot[gray!60] coordinates {(12.80,0) (14.20,2000)};
\addplot[gray]  coordinates {(12.91,0) (14.31,2000)} node[pos=0.55, left, black, font=\scriptsize]{$+20$};
\addplot[gray!60] coordinates {(13.02,0) (14.42,2000)};
\addplot[gray]  coordinates {(13.13,0) (14.53,2000)} node[pos=0.42, left, black, font=\scriptsize]{$+30$};
\addplot[gray!60] coordinates {(13.24,0) (14.64,2000)};
\addplot[gray]  coordinates {(13.35,0) (14.75,2000)} node[pos=0.3, left, black, font=\scriptsize]{$+40$};
\addplot[black, very thick] coordinates {(12.82,0) (13.93,2000)}
  node[pos=0.75, right=2pt, font=\scriptsize]{ISA};
\end{axis}
\end{tikzpicture}
\end{center}

Loi lue sur la planche : pas $= 12{,}03 + 0{,}022\,(T + 20) + 0{,}70\,h/1000$,
avec $T$ en \deg{}C et $h$ en mètres ; en atmosphère normale,
pas $= 12{,}82 + 0{,}55\,h/1000$. Traits intermédiaires tous les 5 \deg{}C.
Lecture à $\pm 0{,}05$ \deg{}.

\section{Masse maximale au décollage avec effet de sol (page 9)}

Départ de terrain dégagé, patins ou flotteurs à 1,5 m du sol,
34\,000 tr/min.

\begin{center}
\begin{tikzpicture}
\begin{axis}[planche, height=10.5cm,
  xlabel={Température extérieure (\deg{}C)}, ylabel={Altitude-pression (m)},
  xmin=-40, xmax=46, ymin=0, ymax=4800,
  minor x tick num=1, minor y tick num=4]
\addplot[black, smooth] coordinates
  {(-12.0,4620) (-5.7,4430) (-2.5,4336) (0.6,4241) (7.0,4054) (13.3,3868) (16.5,3775) (19.7,3683) (26.0,3499)}
  node[pos=0.45, above right, font=\scriptsize]{1100 kg};
\addplot[black, smooth] coordinates
  {(-34.2,4620) (-23.4,4308) (-18.0,4151) (-12.6,3993) (-1.8,3676) (8.9,3356) (14.3,3195) (19.7,3033) (30.5,2708)}
  node[pos=0.55, above right, font=\scriptsize]{1200};
\addplot[black, smooth] coordinates
  {(-40.0,4152) (-25.8,3740) (-18.8,3534) (-11.7,3329) (2.5,2919) (16.7,2509) (23.8,2305) (30.8,2101) (45.0,1694)}
  node[pos=0.55, above right, font=\scriptsize]{1300};
\addplot[black, smooth] coordinates
  {(-40.0,3563) (-25.8,3138) (-18.8,2926) (-11.7,2714) (2.5,2291) (16.7,1869) (23.8,1658) (30.8,1448) (45.0,1028)}
  node[pos=0.55, above right, font=\scriptsize]{1400};
\addplot[black, smooth] coordinates
  {(-40.0,3013) (-25.8,2589) (-18.8,2377) (-11.7,2165) (2.5,1741) (16.7,1319) (23.8,1108) (30.8,897) (45.0,476)}
  node[pos=0.55, above right, font=\scriptsize]{1500};
\addplot[black, smooth] coordinates
  {(-40.0,2490) (-26.3,2072) (-19.5,1863) (-12.6,1655) (1.1,1240) (14.8,825) (21.6,619) (28.4,412) (42.1,0)}
  node[pos=0.55, above right, font=\scriptsize]{1600};
\addplot[black, very thick] coordinates {(15,0) (-16.2,4800)}
  node[pos=0.83, right=2pt, font=\scriptsize]{ISA};
\end{axis}
\end{tikzpicture}
\end{center}

\section{Masse maximale au décollage hors effet de sol (page 10)}

Décollage sous forte pente, ou travail en stationnaire hors effet de sol,
34\,000 tr/min.

\begin{center}
\begin{tikzpicture}
\begin{axis}[planche, height=10.5cm,
  xlabel={Température extérieure (\deg{}C)}, ylabel={Altitude-pression (m)},
  xmin=-40, xmax=46, ymin=0, ymax=4800,
  minor x tick num=1, minor y tick num=4]
\addplot[black, smooth] coordinates
  {(-8.1,4600) (-2.5,4457) (0.3,4385) (3.1,4314) (8.6,4170) (14.2,4026) (17.0,3954) (19.7,3882) (25.3,3737)}
  node[pos=0.45, above right, font=\scriptsize]{1000 kg};
\addplot[black, smooth] coordinates
  {(-30.9,4600) (-20.3,4312) (-15.0,4168) (-9.8,4024) (0.8,3737) (11.4,3451) (16.6,3307) (21.9,3164) (32.5,2878)}
  node[pos=0.55, above right, font=\scriptsize]{1100};
\addplot[black, smooth] coordinates
  {(-40.0,4213) (-25.8,3800) (-18.8,3595) (-11.7,3390) (2.5,2984) (16.7,2582) (23.8,2382) (30.8,2183) (45.0,1788)}
  node[pos=0.55, above right, font=\scriptsize]{1200};
\addplot[black, smooth] coordinates
  {(-40.0,3565) (-25.8,3147) (-18.8,2939) (-11.7,2732) (2.5,2318) (16.7,1908) (23.8,1704) (30.8,1500) (45.0,1094)}
  node[pos=0.55, above right, font=\scriptsize]{1300};
\addplot[black, smooth] coordinates
  {(-40.0,3039) (-25.8,2608) (-18.8,2394) (-11.7,2181) (2.5,1757) (16.7,1337) (23.8,1128) (30.8,920) (45.0,507)}
  node[pos=0.55, above right, font=\scriptsize]{1400};
\addplot[black, smooth] coordinates
  {(-40.0,2502) (-25.9,2075) (-18.9,1864) (-11.9,1653) (2.2,1234) (16.3,819) (23.3,613) (30.3,407) (44.4,0)}
  node[pos=0.55, above right, font=\scriptsize]{1500};
\addplot[black, smooth] coordinates
  {(-40.0,1984) (-28.9,1646) (-23.3,1478) (-17.8,1311) (-6.7,979) (4.4,650) (10.0,486) (15.6,324) (26.7,0)}
  node[pos=0.5, above right, font=\scriptsize]{1600};
\addplot[black, very thick] coordinates {(15,0) (-16.2,4800)}
  node[pos=0.83, right=2pt, font=\scriptsize]{ISA};
\end{axis}
\end{tikzpicture}
\end{center}

Sur les deux planches, les courbes sont tracées sur leur étendue imprimée
réelle : les courbes 1100 et 1200 de la page 9 et la courbe 1000 de la
page 10 s'arrêtent avant le bord du cadre sur l'original.

\subsection*{Plafonds de stationnaire en atmosphère normale}

Croisement de chaque courbe de masse avec la droite ISA, lu à $\pm 50$ m.
Les valeurs marquées d'un astérisque correspondent à un croisement situé
hors du cadre de la planche (extrapolation, $\pm 150$ à $\pm 200$ m) ; ils
ne sont volontairement pas tracés sur les figures.

\begin{center}
\small
\begin{tabular}{@{}rrr@{}}
\toprule
Masse & Avec effet de sol (1,5 m) & Hors effet de sol \\
\midrule
1000 kg & & $\sim$4\,820 m$^{*}$ \\
1100 kg & $\sim$4\,700 m$^{*}$ & 4\,070 m \\
1200 kg & 3\,930 m & 3\,230 m \\
1300 kg & 3\,150 m & 2\,410 m \\
1400 kg & 2\,380 m & 1\,720 m \\
1500 kg & 1\,700 m & 1\,060 m \\
1600 kg & 1\,020 m & 420 m \\
\bottomrule
\end{tabular}

{\scriptsize $^{*}$ extrapolé au-delà du cadre de la planche.}
\end{center}

Pente locale des courbes au voisinage du croisement : $-28$ à $-30$ m par
\deg{}C selon la masse. L'effet de sol à 1,5 m équivaut à environ 85 à
95 kg de masse en moins. En prolongeant l'espacement des croisements, la
masse maximale au stationnaire ISA au niveau de la mer sort vers 1\,665 kg
hors effet de sol et vers 1\,750 kg avec.

\end{document}
